\documentclass[12pt, letterpaper]{article}
\date{\today}
\usepackage[margin=1in]{geometry}
\usepackage{amsmath}
\usepackage{hyperref}
\usepackage{cancel}
\usepackage{amssymb}
\usepackage{fancyhdr}
\usepackage{pgfplots}
\usepackage{booktabs}
\usepackage{pifont}
\usepackage{amsthm,latexsym,amsfonts,graphicx,epsfig,comment}
\pgfplotsset{compat=1.16}
\usepackage{xcolor}
\usepackage{tikz}
\usetikzlibrary{shapes.geometric}
\usetikzlibrary{arrows.meta,arrows}
\newcommand{\Z}{\mathbb{Z}}
\newcommand{\N}{\mathbb{N}}
\newcommand{\R}{\mathbb{R}}
\newcommand{\Q}{\mathbb{Q}}
\newcommand{\C}{\mathbb{C}}
\newcommand{\F}{\mathbb{F}}

\newcommand{\Po}{\mathcal{P}}
\newcommand{\Pro}{\mathbb{P}}
\author{Alex Valentino}
\title{501 homework}
\pagestyle{fancy}
\renewcommand{\headrulewidth}{0pt}
\renewcommand{\footrulewidth}{0pt}
\fancyhf{}
\rhead{
	Homework 10\\
	501
}
\lhead{
	Alex Valentino\\
}
\begin{document}
\begin{enumerate}
	\item[6]
	\begin{align*}
	&\frac{1}{4}(\|f+g\|^2 - \|f - g\|^2 - i \|f+ig\|^2 + i\|f-ig\|^2)\\
	&\frac{1}{4}(\langle f + g, f+ g\rangle - \langle f - g, f - g\rangle - i\langle f + ig, f + ig\rangle + i \langle f - ig, f - ig \rangle)\\
	&\frac{1}{4}(\|f\|^2 + \|g\|^2 + 2 Re(\langle f,g \rangle) - \|f\|^2 - \|g\|^2  + 2 Re(\langle f,g\rangle) +\\
	&i(-\|f\|^2 + \|g\|^2  + 2Im(\langle f,g\rangle) + \|f\|^2 - \|g\|^2 +  2Im(\langle f, g \rangle))\\
	& \frac{1}{4}(4Re(\langle f,g\rangle) + 4 Im (\langle f,g\rangle)\\
	&= \langle f,g\rangle
	\end{align*}
	\item[7] Let $\mathcal{H} = L^2([0,\infty), \mathcal{B}, \mu)$.  We claim that the function
	$f: (0,\infty) \to \mathcal{H}$ given by 
	$f(t) = 1_{[0,t)}$ is continuous and has 
	the property for all $t \in (0,\infty)$
	that for all $r < s < t$ we have that 
	$\langle f(t) - f(s), f(s) - f(r) \rangle$.  Note that $f$ is continuous since 
	for a given $\epsilon > 0$ we have that if $|t-s| < \epsilon^2$ with WLOG 
	$t > s$ then 
	$$
	\|f(t) - f(s)\| = \sqrt{\int_{(0,\infty)} | |1_{(0,t]} - 1_{(0,s]}|^2d\mu} =
	\sqrt{\int_{(0,\infty)} | |1_{(s,t]}|^2d\mu} <
	\sqrt{\epsilon^2} = \epsilon.
	$$
	Furthmore the self orthogonality comes from the fact that 
	\begin{align*}
		\langle f(t) - f(s), f(s) - f(r) \rangle &= \langle f(t), f(s) \rangle + 
		\langle f(s), f(r) \rangle - \langle f(s), f(s) \rangle - \langle f(t), f(r) \rangle \\
		&= s + r - s - r\\
		&= 0
	\end{align*}
	For a general Hilbert space we have that:
		\item[8]
		

	\item[9] Let $\{f_n\}_{n \in \N}$ be a weakly convergent sequenece in a Hilbert 
	space $\mathcal{H}$ with limit $f$.  We want to show that 
	$\lim \| f - f_n\| = 0$ iff $\lim \|f_n\| = \|f\|$.
	\begin{itemize}
		\item $\Rightarrow$ Suppose that $\lim \|f - f_n\| = 0$.  Then we have 
		that $$\lim \|f\|^2 + \|f_n\|^2 - 2 Re(\langle f,f_n \rangle) = 0$$
		Note by the weak convergence of $f_n$ we have that 
		$\lim \|f\|^2 + \|f_n\|^2 = 2\|f\|^2$, therefore 
		$\lim \|f_n\|^2 = \|f\|^2$, which implies our desired result
		\item $\Leftarrow$ Suppose that $\lim \|f_n\| = \|f\|$.  
		Then we have be the parallelogram law that 
		\begin{align*}
		\lim \|f - f_n\|^2 &= \lim 2\|f\|^2 + 2\|f_n\|^2 - \|f + f_n\|^2\\
		&= \lim 4 \|f\|^2 - \|f\|^2 - \|f_n\|^2 - 2Re <f,f_n>\\
		&= 4 \|f\|^2 - \|f\|^2 - \|f\|^2 - 2\|f\|^2\\
		&= 0
		\end{align*}
		Note that the last line is true because of the weak convergence of $f_n$.
		Therefore $\lim \|f - f_n\| = 0$.
		\iffalse Then by the 
		identity proven in the proof of cauchy-schwarz we have that
		$\|f\|^2 = \lim |\langle f, f_n\rangle| = \lim \|f\| \|f_n\| - \frac{1}{2}\|f\|\|f_n\|\|f/\|f\| - f_n/\|f_n\|\|^2$ which implies that 
		$0 = \lim \|f\|\|f_n\|\|f/\|f\| - f_n/\|f_n\|\|^2 = \lim \|f - f_n\|^2$.
		Note that one can exclude the scenario where $\|f\| = 0$ since if it was 
		then $f= 0$ and the convergence in norm is trivially true.  Therefore
		\fi
	\end{itemize}
	\item[10]
	Let $\mu$ be the Borel measure on $[0,1]$, $C$ be the set of real valued measurable
	functions $f$ on $[0,1]$ such that $\int_{[-1,1]} f^2 d \mu = 1, 
	\int_{[-1,1]} x^j f d \mu = 0$ for $j \in \{0,1,2\}$.  
	The maximizing value of $\int_{[-1,1]} x^3 f(x) d\mu$ is 
	$\frac{2\sqrt{2}}{5\sqrt{7}}$, however I don't know how to prove this.  
	I believe if I applied gram-schmidt then there would've been a point in which 
	the inner product was taken between computable polynomials, however I 
	did not have enough time to do this.
\end{enumerate}
\end{document}
